\section{Methods}

\subsection{Participants and data}
We used the 128-channel event-related potential (ERP) recordings of MODMA (Lanzhou, 2014--2016)~\cite{cai2022modma}.
The release contains 53 recordings, one per participant, sampled at \SI{250}{Hz}. They were acquired with a HydroCel
Geodesic Sensor Net (EGI) using 128 scalp electrodes plus the vertex reference (E129/Cz).

Diagnostic group (24 MDD, 29 HC) was taken \emph{exclusively} from the official subject-information spreadsheet.
Recording files were matched to it by exact identifier, without normalisation: all 53 identifiers matched one-to-one.

Table~\ref{tab:demo} summarises the sample. The groups did not differ in age ($p=\DemoAgeP$) or sex
($p=\DemoGenderP$), but MDD participants had fewer years of education ($p=\DemoEducationP$). Questionnaire scores
(PHQ-9, GAD-7, CTQ-SF, LES, SSRS, PSQI) were used only to describe the sample and never as predictors, because they
measure, or are closely tied to, the target.

\begin{table}[t]
\centering\small
\caption{Sample characteristics (all 53 recordings). Mean $\pm$ SD. $^\dagger$Descriptive only; never used as
predictors.}
\label{tab:demo}
\begin{tabular}{@{}lllll@{}}
\toprule
Variable & MDD ($n=24$) & HC ($n=29$) & Test & $p$ \\
\midrule
Age (years) & 30.9 $\pm$ 10.4 & 31.4 $\pm$ 9.1 & Mann--Whitney & 0.82 \\
Education (years) & 13.2 $\pm$ 3.8 & 16.0 $\pm$ 2.3 & Mann--Whitney & 0.0048 \\
PHQ-9$^\dagger$ & 18.3 $\pm$ 3.5 & 2.7 $\pm$ 1.8 & Mann--Whitney & $<$0.001 \\
GAD-7$^\dagger$ & 13.4 $\pm$ 4.9 & 2.1 $\pm$ 2.1 & Mann--Whitney & $<$0.001 \\
PSQI$^\dagger$ & 12.5 $\pm$ 4.2 & 4.1 $\pm$ 2.2 & Mann--Whitney & $<$0.001 \\
Sex (F/M) & 11/13 & 9/20 & Fisher exact & 0.39 \\
\bottomrule
\end{tabular}
\end{table}

\paragraph{Task and events.}
Each trial of the dot-probe task~\cite{macleod1986attentional} consisted of the following sequence:
\begin{enumerate}
  \item a fixation cross (464--652~ms);
  \item a pair of faces, one emotional and one neutral (516~ms);
  \item an inter-stimulus interval (116--316~ms);
  \item a dot probe (168~ms), followed by the participant's response.
\end{enumerate}
Trials came in three blocks of 160 (fear, happy, sad), with block order varying across participants.

We decoded events from the individual EGI event channels by name (\texttt{fcue}, \texttt{hcue}, \texttt{scue}). The
combined stimulus channel produced by the reader carries the same code for every event in these files and is
therefore unusable. Every recording contained the complete 480 trials. Epochs were locked to face-pair onset, because
the probe appears at least 632~ms after it, so the window up to 600~ms is free of probe-evoked activity.

\subsection{Pre-specified analysis plan}
Before any epoch, feature or model output existed, we committed an analysis plan to version control. It fixed the
preprocessing, the participant exclusion criteria, the feature definitions, the evaluation protocol, the primary model
and the primary outcome.

Five deviations are listed in Appendix~\ref{app:deviations}: four were decided before any model was run, and the
fifth (a computational shortcut for the Riemannian models) before any of their results existed. The only one
affecting the data was replacing an ad-hoc bad-channel rule with the PREP implementation (below), after the original
rule produced implausible flags next to the reference electrode.

\subsection{Preprocessing}
All steps were performed separately for each participant with MNE-Python~\cite{gramfort2013mne}, so no information
is shared across participants:
\begin{enumerate}
  \item Bad channels were detected with the PREP pipeline (PyPREP \texttt{NoisyChannels}, without
        RANSAC)~\cite{bigdely2015prep}. The criteria were NaN/flat signal, deviation, high-frequency noise, low
        correlation, low SNR, dropout and abnormal PSD. The vertex reference was exempt, being flat by construction.
  \item Independent component analysis (Picard, extended-infomax objective~\cite{ablin2018picard}, 30 components)
        was fitted on a copy filtered at 1--100~Hz, notched at 50~Hz (the mains frequency in China) and
        average-referenced, with bad channels excluded.
  \item Components were labelled with ICLabel~\cite{pion2019iclabel,li2022mneicalabel}. A component was removed when
        its most probable class was eye, muscle, heart, line noise or channel noise \emph{and} that class had
        probability $\ge 0.80$.
  \item The analysis copy was band-pass filtered at 0.5--40~Hz (zero-phase FIR) and notched at 50~Hz. It was then
        average-referenced, the ICA cleaning was applied, bad channels were interpolated (spherical splines), and it
        was re-referenced to the average.
  \item Epochs spanned $-200$ to $600$~ms, with a $-200$ to $0$~ms baseline. An epoch was rejected if any channel
        exceeded $\pm\SI{100}{\micro\volt}$.
  \item For spectral, connectivity and complexity features, the task period of the same cleaned recording was
        resampled to 125~Hz and cut into 2-s segments. The same amplitude rule was applied to the segments.
\end{enumerate}

\paragraph{Participant exclusion (pre-specified).}
A participant was excluded if any of these held:
\begin{itemize}
  \item more than 20\% of channels were bad;
  \item fewer than 240 of 480 trials (50\%) survived rejection;
  \item fewer than 40 trials survived in any condition.
\end{itemize}

\subsection{Features}
All features were computed per participant, without labels. They are therefore not a source of leakage; everything
that \emph{learns} from data happens inside the cross-validation loop.

\paragraph{Regions.}
We defined ten regions: frontal, central, temporal, parietal and occipital, each split into left and right
(Fig.~\ref{fig:regions} in the appendix). Each HydroCel electrode was assigned to the region of its nearest 10--10
anchor position on the fitted head sphere. Midline electrodes and the lower face/neck ring (more than \SI{115}{\degree}
from the vertex) were not used for regional features.

\paragraph{Feature families.}
\begin{itemize}
  \item \textbf{Spectral (100 features).} Welch power in the delta (1--4~Hz), theta (4--8), alpha (8--13),
        beta (13--30) and low-gamma (30--40~Hz) bands, as $\log_{10}$ absolute and relative power per region.
  \item \textbf{Frontal alpha asymmetry (2).} $\ln\alpha_\text{right}-\ln\alpha_\text{left}$ for the HydroCel
        equivalents of F4/F3 (E124/E24) and F8/F7 (E122/E33)~\cite{allen2015faa,vandervinne2017faa}.
  \item \textbf{ERP (21).} Measured per condition on the participant's average, using windows fixed a priori from
        the literature:
        \begin{itemize}
          \item N170 peak amplitude and latency at left and right occipito-temporal clusters (around P7/P8),
                140--200~ms~\cite{bentin1996n170};
          \item P300 mean amplitude (250--500~ms) and peak latency at a centro-parietal cluster
                (around Pz)~\cite{polich2007p300};
          \item late positive potential (LPP), mean amplitude 400--600~ms, at the same
                cluster~\cite{hajcak2010erp}.
        \end{itemize}
  \item \textbf{Connectivity (135).} Weighted phase-lag index (wPLI)~\cite{vinck2011wpli} between the ten
        region-mean signals, in the theta, alpha and beta bands.
  \item \textbf{Complexity (30).} Per region, sample entropy~\cite{richman2000sampen}, Higuchi fractal
        dimension~\cite{higuchi1988} and Lempel--Ziv complexity~\cite{lempel1976complexity}, averaged over 2-s
        segments.
\end{itemize}
This gives 288 subject-level features. For trial-level models we kept the ten region-mean signals of each epoch
($10\times201$ samples) and, for deep learning, the full $129\times201$ epochs.

\subsection{Models}
\begin{itemize}
  \item \textbf{Demographic baseline.} L2 logistic regression on age, sex and education. It indicates how much of
        any EEG result could be explained by demographics alone.
  \item \textbf{Handcrafted features.} Each model is a scikit-learn~\cite{pedregosa2011sklearn} pipeline: median
        imputation, standardisation, then the classifier.
        \begin{itemize}
          \item \emph{Primary, fixed a priori:} L2-regularised logistic regression (LR-L2) on all 288 features.
          \item \emph{Secondary:} L1-LR, linear and RBF support-vector machines, random forest~\cite{breiman2001rf}
                and depth-limited LightGBM~\cite{ke2017lightgbm}, plus LR-L2 on each feature family separately.
        \end{itemize}
        Picking the best of many models inflates its estimate~\cite{cawley2010overfitting}. We therefore also report a
        \emph{nested model selection} estimate, in which the inner loop chooses both the model family and its
        hyperparameters.
  \item \textbf{Riemannian, trial level.}~\cite{barachant2012multiclass,congedo2017riemannian}
        \begin{itemize}
          \item OAS-shrunk covariance matrices~\cite{chen2010oas} of the regional epoch signals, either projected
                to the tangent space and classified with LR, or classified with minimum distance to mean (MDM);
          \item ERP-augmented covariances, whose class prototypes are estimated inside each training fold, followed
                by tangent space and LR with a fixed $C=0.1$ (deviation D5).
        \end{itemize}
  \item \textbf{Deep learning, trial level.} EEGNet~\cite{lawhern2018eegnet}, ShallowFBCSPNet~\cite{schirrmeister2017deep}
        and a temporal convolutional network (TCN)~\cite{bai2018tcn,gemein2020machine}, as implemented in
        braindecode.
        \begin{itemize}
          \item Training used AdamW (learning rate $10^{-3}$, weight decay $10^{-3}$), cosine annealing,
                class-weighted cross-entropy, and augmentation by random time shifts ($\pm40$~ms), amplitude
                scaling and Gaussian noise.
          \item Early stopping (patience 10, maximum 60 epochs) used the subject-level AUC on a validation split of
                the \emph{training participants} (20\%).
          \item Channel-wise normalisation statistics came from the inner training participants only.
        \end{itemize}
\end{itemize}
Trial-level models produce one probability per trial. These were averaged within each participant to give a single
participant-level prediction.

\subsection{Evaluation protocol}
\begin{itemize}
  \item \textbf{Outer loop.} Stratified 5-fold cross-validation over \emph{participants}, repeated 10 times
        (3 times for deep learning, for computational reasons). All models were evaluated on \emph{identical}
        participant partitions.
  \item \textbf{Inner loop.} Hyperparameters were selected by 5-fold participant-grouped cross-validation within the
        training participants, maximising AUC.
  \item \textbf{Leakage check.} An assertion verifies, in every outer, inner and validation split, that no
        participant contributes data to both sides.
  \item \textbf{Metrics.} All metrics are computed per participant: ROC AUC (the pre-specified primary outcome),
        balanced accuracy, sensitivity and specificity at a fixed 0.5 threshold (never tuned on test data), F1 and
        Brier score.
  \item \textbf{Uncertainty.} Point estimates are averaged over repetitions. 95\% confidence intervals come from
        2{,}000 stratified bootstrap resamples of \emph{participants} applied to the out-of-fold predictions, with
        interval bounds averaged over repetitions.
  \item \textbf{Significance.} Label-permutation tests~\cite{combrisson2015exceeding} permute labels \emph{between
        participants} and re-run the entire nested pipeline. Fast models used 1{,}000 permutations; slower models
        used 200. We report $p=(1+\#\{\text{AUC}_\text{perm}\ge\text{AUC}_\text{obs}\})/(1+n_\text{perm})$.
\end{itemize}

\subsection{Confounds, sensitivity and inflation control}
\begin{itemize}
  \item \textbf{Can EEG predict confounds?} We applied the primary pipeline to three alternative targets: sex,
        age above the median, and recording period.
        \begin{itemize}
          \item The recording period is relevant because group and period are partly confounded: between October 2015
                and January 2016 only controls were recorded, and in March--April 2016 only patients.
          \item We therefore also re-ran the primary pipeline on the participants recorded between June and
                August 2015, when both groups were recorded in parallel.
        \end{itemize}
  \item \textbf{Exclusions.} Because exclusions were more frequent in MDD, we re-ran the primary model on all 53
        participants.
  \item \textbf{Signal quality.} We compared retained trials, interpolated channels and removed ICA components
        between groups.
  \item \textbf{Inflation control (invalid by design).} To quantify the inflation caused by trial-wise splitting,
        we evaluated two trial-level models with both a trial-wise split and a subject-wise split, keeping
        everything else identical:
        \begin{itemize}
          \item the tangent-space model;
          \item a random forest on raw regional amplitudes, mimicking common practice.
        \end{itemize}
\end{itemize}

\subsection{Interpretation}
\begin{itemize}
  \item \textbf{Permutation importance.} For the primary model, computed on the \emph{test} folds of the outer loop,
        for single features and for whole feature families permuted jointly.
  \item \textbf{SHAP.} SHAP values~\cite{lundberg2017shap} of the primary model refitted on all participants
        (descriptive).
  \item \textbf{Univariate effects.} Cohen's $d$ with Mann--Whitney tests and Benjamini--Hochberg false-discovery-rate
        correction~\cite{benjamini1995fdr}.
  \item \textbf{Cluster-based permutation tests}~\cite{maris2007nonparametric} on channel-level band power
        (spatial clusters) and on ERPs (spatio-temporal clusters), as descriptive neurophysiology.
\end{itemize}
